\begin{abstract}
Pose-based gait data are widely used for clinical screening, yet whether Large Language Models (LLMs) can interpret such motion evidence remains underexplored.
We introduce GaitEval, a six-dataset benchmark for evidence-grounded label prediction across ASD-Kinect, ATAXIA, CARE-PD, GMFCS, PARKINSONISM, and PATHOGAIT, with heterogeneous populations, pose formats, and label spaces.
A baseline-anchored protocol compares controlled pose-derived inputs with label-prior and same-feature linear references, and evaluates zero-shot, few-shot, and low-rank adaptation (LoRA) within their tested coverage.
The retained five-family motion-summary evaluation shows near-chance agreement on its original tasks; a shared six-dataset, subject-disjoint Qwen2.5-7B evaluation identifies important dataset--representation exceptions.
Linear references recover task-relevant signal where available, while adaptation recovers some in-distribution mappings: standardized GMFCS summary macro-F1 rises from 0.265 to 0.467, but remains below its matched linear reference of 0.579.
The accepted-protocol ATAXIA result likewise retains its agreement increase from $\kappa=0.046$ to $0.507$; transfer remains unreliable in the tested bidirectional ATAXIA--PATHOGAIT probe.
Overall, evidence use is heterogeneous and representation-dependent: adaptation can learn dataset-specific mappings, while reliability beyond the tested settings remains unresolved.
\end{abstract}
